% !TeX root = ../../../../../main.tex
% chapters/sections/chapter5/subsections/two_plane/overview.tex

\subsection{2平面分析のための4曲線の導出の概要}
\label{sec:chapter5_two_plane_overview}
本節ではシミュレーション結果を解釈するための理論的枠組みとして、
\ref{chap:chapter3}
章の方程式系から自国の
\textbf{YS（生産供給）曲線}、
\textbf{YD（生産需要）曲線}、
\textbf{CS（消費供給）曲線}および
\textbf{CD（消費需要）曲線}
を導出する。
YS曲線およびCS曲線は主に自国代表的家計の供給側の式を集約したものであり、
YD曲線およびCD曲線は主に自国代表的家計の需要側の式を集約したものである。
そして
YS曲線およびYD曲線を用いることにより\((y_t^H,p_t^H)\)平面上の均衡点の動きを、
CS曲線およびCD曲線を用いることにより\((c_t^{H\to W},p_t^{H\to W})\)平面上の均衡点の動きを
視覚的に理解することが可能になる。